\documentclass[tikz,border=3pt]{standalone}
\input{tikz-tcee}
\begin{document}
\hyphenpenalty=10000 \exhyphenpenalty=10000
% Semáforo de los umbrales oficiales: verde dentro, rojo fuera (alerta)
% Barras primero, después las marcas del umbral y, al final, los rótulos.
\begin{tikzpicture}[tcee, font=\scriptsize,
  ok/.style={fill=tceeverde!70}, mal/.style={fill=tceerojo!75},
  n/.style={anchor=east, font=\scriptsize, text width=2.75cm, align=right},
  u/.style={font=\scriptsize\bfseries, fill=white, inner sep=1pt}]
  % barras (0 a 5 cm cada una)
  \fill[mal] (0,0) rectangle (1.2,0.32); \fill[ok] (1.2,0) rectangle (3.8,0.32); \fill[mal] (3.8,0) rectangle (5,0.32);
  \fill[mal] (0,-0.6) rectangle (2.0,-0.28); \fill[ok] (2.0,-0.6) rectangle (5,-0.28);
  \fill[ok] (0,-1.2) rectangle (2.5,-0.88); \fill[mal] (2.5,-1.2) rectangle (5,-0.88);
  \fill[ok] (0,-1.8) rectangle (2.5,-1.48); \fill[mal] (2.5,-1.8) rectangle (5,-1.48);
  \fill[ok] (0,-2.4) rectangle (2.5,-2.08); \fill[mal] (2.5,-2.4) rectangle (5,-2.08);
  \fill[ok] (0,-3.0) rectangle (2.5,-2.68); \fill[mal] (2.5,-3.0) rectangle (5,-2.68);
  \fill[ok] (0,-3.6) rectangle (2.5,-3.28); \fill[mal] (2.5,-3.6) rectangle (5,-3.28);
  \fill[ok] (0,-4.2) rectangle (2.5,-3.88); \fill[mal] (2.5,-4.2) rectangle (5,-3.88);
  % marcas de los umbrales
  \foreach \x/\y in {1.2/0, 3.8/0, 2.0/-0.6, 2.5/-1.2, 2.5/-1.8, 2.5/-2.4, 2.5/-3.0, 2.5/-3.6, 2.5/-4.2}
    \draw[very thick, dicinflacion] (\x,\y-0.05) -- (\x,\y+0.37);
  % rótulos
  \node[n] at (-0.1,0.16) {Cuenta corriente (media 3 años)};
  \node[n] at (-0.1,-0.44) {PIIN};
  \node[n] at (-0.1,-1.04) {Paro};
  \node[n] at (-0.1,-1.64) {CLU nominales (3 años, euro)};
  \node[n] at (-0.1,-2.24) {Precio de la vivienda (anual)};
  \node[n] at (-0.1,-2.84) {Deuda de los hogares};
  \node[n] at (-0.1,-3.44) {Deuda de las empresas};
  \node[n] at (-0.1,-4.04) {Déficit / deuda públicos};
  \node[u] at (1.2,0.16) {$-4$}; \node[u] at (3.8,0.16) {$+6$};
  \node[u] at (2.0,-0.44) {$-35$};
  \node[u] at (2.5,-1.04) {10};
  \node[u] at (2.5,-1.64) {$+9$};
  \node[u] at (2.5,-2.24) {$+9$};
  \node[u] at (2.5,-2.84) {55};
  \node[u] at (2.5,-3.44) {85};
  \node[u] at (2.5,-4.04) {3 / 60};
  \node[anchor=north west, text width=7.9cm, font=\scriptsize\itshape, text=tceegris] at (-2.85,-4.35)
    {En \% (del PIB, salvo paro y variaciones). Verde: dentro del umbral; rojo: alerta del procedimiento de desequilibrios (PDM) o del Tratado.};
\end{tikzpicture}
\end{document}
